\begin{pastexam}[label={px:ncu-113-dsa-02}]{中央資工所 113 學年度 第 2 題}
（複選；每題 5 分，全部答對才給分，答錯倒扣一分）

Consider the following five sorting algorithms: Quick Sort, Heap Sort, Merge Sort, Insertion Sort and Radix
Sort. Which of the following statements are true?
\begin{enumerate}[label=(\Alph*)]
\item Insertion Sort gives best performance when applied on an array which is sorted or almost sorted
(maximum 1 or two elements are misplaced).
\item Radix Sort is most efficient to sort string consisting of ASCII characters?
\item All these five sorting algorithms are comparison based sorting algorithms.
\item All these five sorting algorithms are stable.
\end{enumerate}
\end{pastexam}

\begin{pastsoln}
答案：(A)(B)。

(B) 句尾多了一個問號，照敘述句來判斷；「most efficient」是和題目列的五種排序相比。

\begin{itemize}
\item (A) 對。 插入排序的時間是 $\Theta(n + I)$，$I$ 是逆序對個數（資結書第~14.1~節）。 已經排好時
$I = 0$，只比 $n - 1$ 次；只有一兩個元素放錯位置時，$I$ 最多是 $O(n)$，時間仍是 $\Theta(n)$。 其他四種
在這種輸入上沒有更快的：合併排序、堆積排序還是 $\Theta(n\lg n)$；快速排序固定取頭或尾當基準時，排好的
輸入反而是最壞情況 $\Theta(n^2)$；基數排序不管輸入順序，固定做 $d$ 輪。
\item (B) 對。 ASCII 字元只有 $128$ 種，可以把每個字元直接當成一位數，基數排序只要一輪計數排序
（資結書第~15.2~節），$\Theta(n + 128) = \Theta(n)$。 排一串字串時，每個字元一位，LSD 基數排序
（資結書第~15.2.2~節）是 $\Theta(d(n + 128))$，對輸入的總長度是線性的。 其他四種都是比較排序，
至少要 $\Omega(n\lg n)$ 次比較（資結書定理~15.1），所以五種裡基數排序最快。
\item (C) 錯。 基數排序不比較兩個鍵的大小，它看的是每一位的值，用計數排序分配，所以不受
$\Omega(n\lg n)$ 下界的限制。
\item (D) 錯。 快速排序和堆積排序都不穩定（資結書第~14.7~節的比較表）。 例如堆積排序排
$\langle 2_a, 2_b, 1\rangle$：建堆後 $2_a$ 在根，第一次把根和最後一格的 $1$ 交換，$2_a$ 被放到最後，
輸出是 $1, 2_b, 2_a$，兩個 $2$ 的相對順序反了。
\end{itemize}
\end{pastsoln}
